%document class
\documentclass[10pt,oneside]{book}
%%%% Page Info + Commands %%%%%{

%packages
\usepackage{pgfplots}
\pgfplotsset{compat=1.18}
\usepackage{amsmath}
\usepackage{amsfonts}
\usepackage{amsthm,letltxmacro}
\usepackage{amssymb}
\usepackage{mathtools}
\usepackage{enumerate}
\usepackage{enumitem}
\usepackage{array}
\usepackage[dvipsnames]{xcolor}
\usepackage{changepage}

%This will shrink the page
\newcommand{\smallpage}{  \setlength \oddsidemargin{.25in}
            \setlength \textwidth{5.5in}
            \setlength \topmargin{0in}
            \setlength \textheight{9in}}


% Commands for sets of numbers
\newcommand{\Z}{\mathbb{Z}}\newcommand{\R}{\mathbb{R}}\newcommand{\N}{\mathbb{N}}

% gordie spacing and boxing commands
\newcommand{\gspc}{\vspace{0.13cm}} % 0.5 space
\newcommand{\gline}{\gspc\\} % 1.5 space
\newcommand{\gbline}{\gspc\gspc\\} % double space
\newcommand{\gbox}[1]{\fbox{\parbox{\linewidth}{#1}}} % multiline box command


\newcommand{\gimage}[3]{
\begin{figure}[h]   
    \centering          
    \includegraphics[width=#2\textwidth]{#1}  
    \caption{#3}
    \label{fig:#1}
\end{figure}\\
}

\newcommand{\centerit}[1]{{\begin{center} #1 \end{center}}}
\newcommand{\ul}[1]{\underline{#1}}

%%%%%%%% command for graphics %%%%%%%%%%%%%
\usepackage{fancyhdr}
%}

%%%% Page Setup %%%%%%%
\smallpage\sloppy
\pagestyle{fancy}
\lhead{\large{\textbf{Day 3 - Berliner lectures us hooray}}}
%\rfoot{\thepage/\pageref{LastPage} }
\setlength{\headheight}{10pt}


% Proof writing
\LetLtxMacro\oldproof\proof
\let\endoldproof\endproof
\renewenvironment{proof}[1][\proofname]
  {\vspace{0.2cm}\begin{adjustwidth}{2em}{2em}
   \oldproof[#1]}
  {\endoldproof
   \end{adjustwidth}}



\begin{document}

\newcommand{\cg}[1]{\color{OliveGreen}#1\color{black}}
\newcommand{\cb}[1]{\color{BlueViolet}#1\color{black}}
%
% Introductions
%
\subsection*{Last Time in Class}
Last class, we primarily covered:
\begin{itemize}[itemsep=0pt, parsep=0pt]
	\item Euclidean Algorithm
	\begin{itemize}
		\item to find $\gcd(a,b)$
		\item to get $\gcd(a,b)$ as a combo of $a,b$. 
	\end{itemize}
	\item Solving $ax + by = c$. 
  \begin{itemize}
    \item Existence of solutions
    \item Procedure for finding solutions.
  \end{itemize}
\end{itemize}
\begin{enumerate}
  \item Suppose you know that whenever Dorm Room \#K is occupied, then Dorm Room \#($K+1$) will also be occupied.\\
  \textbf{Group Discussion: }Which dorm rooms are guaranteed to be occupied?\gline
  All dorm rooms with numbers $n\ge K$ will be occupied, but no specific dorm room has any guarantee.\gline
  Suppose you know that:
  \begin{itemize}
    \item Dorm Room \#1 is occupied \textit{and}
    \item whenever Dorm Room \#K is occupied, then Dorm Room \#(K+1) will also be occupied. \\
    \textbf{Group Chit-Chat: }Which dorm rooms are guarenteed to be occupied?\gline
    All dorm rooms with number $n \ge 1$.
  \end{itemize}
  \item Suppose you know that:
  \begin{itemize}
    \item Dorm Room \#2 and \#3 are occupied \textit{and}
    \item Whenever Dorm Room \#K is occupied, then Dorm Room \#($K+3$) will also be occupied.\gline
    \textbf{Group Dialogue:} Which dorm rooms are guaranteed to be occupied?\gline
    All dorm rooms that do not match the pattern $1 + 3n, \;n\in\Z \ge 0$.  
  \end{itemize}
  \item Suppoe you know that:
  \begin{itemize}
    \item Dorm Room \#2 is occupied \textit{and}
    \item whenever Dorm Room \#K is occupied, then dorm room $\#(2k)$ will also be occupied. \gline
    \textbf{Group Schmooze:} Which dorm rooms are guarenteed to be occupied?\gline
    All dorm room such that $2^n,\; n\ge 1$. 
  \end{itemize}
\end{enumerate}

\subsection*{Key Takeaways from Warmup}
Here's the main takeaways from the lesson today:
\[
\begin{rcases}
  \text{Suppose Statement \#1 is \cg{true}.}\\
  \text{If statement $K$ is true, then statement $K+1$ is true.\;}
\end{rcases} \text{implies all statements are true}
\]
This is the Principle of (Weak) Mathematical Induction. 
\pagebreak
\subsection*{Practice}
Now, consider the following claim:\gspc
\begin{adjustwidth}{2em}{2em}
  \textit{Claim.} The sum of the first $n$ (positive) integers is $\frac{n(n+1)}2$
\end{adjustwidth}\gspc
Each value of $n$ is its own ``claim'' that needs verification. \gline
Gauss explains this because we can ``pair'' up each number. Consider:
\begin{align*}
  1, 2, \ldots,n
\end{align*}
If we add the first at the last number, we get $1 + n$. If we add up the second and second to last number, we get $2 + (n-1) = 1+n$. So on and so forth. So, we have a total of $n/2$ pairs of $n+1$, so we have our claim. \gline
Now, trying to do this with a formal proof, we write:
\begin{proof}
  \textit{Proceed using induction.}\gline
  \textbf{Base Case: } Let $n= 1$. The sum of the first integer is 1, and:
  \begin{align*}
    \frac{n(n+1)}2 = \frac{1(2)}2 = 1
  \end{align*}\gline
  Now we have shown the first step to be true.\gline
  \textbf{Inductive Step: } Now that we know that statement $K$ is true, we show that the $K+1$ case is true. ``Assume the claim holds when $n=K$.''\gline
  Since the claim is true for $n = K$, we know for a fact that:
  \begin{align*}
    \frac{K(K+1)}2.
  \end{align*}
  Consider $n = K +1$. The sum of the first $K+1$ integers is:
  \begin{align*}
    1 + 2+ 3+\ldots + K + (K+1)= \frac{K(K+1)}2 + (K+1) \text{ by hypothesis.}
  \end{align*}
  Simplifying gives us:
  \begin{align*}
    \frac{K(K+1)}2 + \frac{2(K+1)}2 &= \frac{(K+2)(K+1)}2=\frac{n(n+1)}2
  \end{align*}
  Wabam wow we did it. By induction the claim is true. 
\end{proof}

\subsection*{Little More Practice}
Consider the claim that for all $n\ge 0$, $5\mid 7^n - 2^n$.
Let's do it!
\begin{proof}
  \textit{By Induction.} We want to show that $\forall n \ge 0$, $5\mid 7^n - 2^n$.\gline
  \textbf{Base Case: }Let $n =0$. $7^n - 2^n = 7^0-2^0 = 1 - 1 0$. And,$5\mid 0$.  \gline
  \textbf{Inductive Step: }Assume that $5\mid 7^k - 2^k$. We want to prove that $5\mid 7^{K+1} - 2^{K+1}$.\gline
  Consider the following step:
  \begin{align*}
    7^{K+1}-2^{K+1}&=-7^K7 - 2^K2\\
    &=7^K(5+2) - 2^K2\\
    &= 7^K(5+2) - 2^K\cdot 2\\
    &= 7^K\cdot 5 - 2(7^K - 2^K)
  \end{align*}
  And because $7^K\cdot 5$ is divisible by 5 and $7^K-2^K$ is, $5\mid 7^n-5^n, n\ge 0$.
\end{proof}

\begin{enumerate}
  \item \begin{enumerate}
    \item What is a formula (in terms of $n$) for the $n$-th positive \textbf{odd} integer?\gline
    \[2n- 1\]
    \item \textbf{Group Chart Fill-in:}
    \begin{center}
      \begin{tabular}{|c|c|}
        \hline
        $n$ & \begin{tabular}{c}\textit{sum} of the first $n$ \\positive \textbf{odd} integers\end{tabular}\\\hline\hline
        1 & 1\\\hline
        2 & 4 \\\hline
        3 & 9 \\\hline
        4 & 16 \\\hline
        5 & 25 \\\hline\hline
      \end{tabular}
    \end{center}
    \item \textbf{Group Conjecture: }What do you think is the sum of the first $n$ positive integers?
    \[n^2\]
    \item What would the ``base case'' be? Try and prove the base case\gline
    \[\text{Sum is 1. }1^2 = 1\]
    What would the ``inductive hypothesis'' be?
    \[n = K\implies K^2 = \sum_{i=0}^K(2i - 1)\] 
    Let's prove your conjecture using induction.
    \begin{align*}
      \sum_{i=0}^{K+1}(2i - 1) &= \sum_{i=0}^K(2i - 1) + (2K + 1)\\
      &= K^2 + (2K + 1)\\
      &= (K+1)^2
    \end{align*}
  \end{enumerate}
  \item \begin{enumerate}
    \item \textbf{Group play-time: } Take a good guess at the values of $n$ for which $n! > 3^n$. \gline
    I'm guessing that this applies for all $n\ge 7$, because we have:
    \begin{gather*}
      3 \cdot 3 \cdot 3 \cdot 3 \cdot 3 \cdot 3 \cdot 3\\
      1 \cdot 2 \cdot 3 \cdot 4 \cdot 5 \cdot 6 \cdot 7
    \end{gather*}
    Which simplifies to:
    \begin{gather*}
      3 \cdot 3 \cdot 3 \cdot 3 \cdot 3 \cdot 3\\
      1 \cdot 2 \cdot 4 \cdot 5 \cdot 6 \cdot 7
    \end{gather*}
    And then we get:
    \begin{align*}
      6 &> 3\\
      1\cdot 2 \cdot 5 &> 3\cdot 3\\
      4\cdot 7 =28&> 3 \cdot 3 \cdot 3 =27
    \end{align*}
    It's at this point that all the multipliers of our factorial outnumber the 3s so we achieve victory. 
    \item Prove your guess using induction.\gline
    We will prove the claim that for $n\ge 7$ that $n! > 3^n$.
    \begin{proof}
      \textit{By Induction. }We want to show via. induction that for $n\ge 7$, $n! > 3^n$.\gline
      \textbf{Base Case: }Let $n= 7$. Thus we have that:
      \begin{align*}
        3^n = 2187 &< n! = 5040
      \end{align*}
      \textbf{Inductive Step: }Assume the claim holds when $n = K$. Let $n = K+1$. Thus we have that:
      \begin{align*}
        3^K < K!
      \end{align*}
      Now consider the $K+1$ case:
      \begin{align*}
        3^{K+1}&=3^K3\\
        (K+1)! &= K! \cdot (K+1)
      \end{align*}
      Because in our base case, $n> 7$, we know that $(K+1) > 3$. Thus, we have:
      \begin{gather*}
        3^K\cdot 3 < K!\cdot 3 < K!\cdot(K+1)\\
        3^{K+1} < (K+1)!\\
        3^n < n!
      \end{gather*}
      Hence, for $n\ge 7$, $3^n < n!$.
    \end{proof}
  \end{enumerate}
  
  \item The Fibonacci numbers are a really cool list (i.e sequence) of integers given recurively by:
  \begin{center}
    \begin{tabular}{c c c c}
      $F_1 = 1$ & $F_2 = 1$ & \text{and} & $F_n = F_{n-1} + F_{n-2}$ (for $n\ge 2$)
    \end{tabular}
  \end{center}  
  \begin{enumerate}
    \item List the first ten Fibonacci numbers
    \begin{center}
    \begin{tabular}{|c|c|c|c|c|c|c|c|c|c|c|}
      \hline
      $n$& 1&2 &3 &4 &5 &6&7&8&9&10\\\hline
      $F_n$&1&1&2&3&5&8&13&21&34&55 \\\hline
    \end{tabular}
  \end{center}
    \item Find the \textit{sum of the squares} of the first $n$ Fibonacci numbers.
    \begin{center}
      \begin{tabular}{|c|c|c|c|c|c|c|c|c|c|}
        \hline
        $n$& 1&2 &3 &4 &5 &6&7&8&9\\\hline
        $F_1^2+\ldots+F_n^2$&1&2&6&15&40&104&273&714&?\\\hline
      \end{tabular}
    \end{center}
    \item Make a conjecture regarding $F_1^2+\ldots+F_n^2$.\gline
    $F_1^2+\ldots+F_n^2$ = $F_n \cdot F_{n+1}$
    \item Use induction to prove your conjecture
    \begin{proof}
      \textit{By Induction. } Let $F_n$ be the $n$-th number in the Fibonacci sequence. We want to show via induction that $F_1^2+\ldots+F_n^2$ = $F_n \cdot F_{n+1}$.\gline
      \textbf{Base Case: }Let $n= 1$. Consider that:
      \begin{align*}
        F_1^2 &= 1^2 = 1\\
        F_1\cdot F_2 &= 1\cdot 1 = 1
      \end{align*}
      \textbf{Inductive Step: } We know that $n= K$ holds. So, we asumme that:
      \begin{align*}
        F_1^2+\ldots + F_K^2 = F_K\cdot F_{K+1}
      \end{align*}
      Now we want to show that the $n = K+1$ case holds. Lookie here:
      \begin{align*}
        F_1^2 + \ldots + F_K^2 + F_{K+1}^2 &= F_{K}\cdot F_{K+1} + F_{K+1}^2\\
        &= F_K\cdot F_{K+1} + F_{K+1}F_{K+1}\\
        &= F_{K+1}(F_K+F_{K+1})\\
        &= F_{K+1}\cdot F_{K+2}
      \end{align*}
      yay we did it we're done. end of proof.
    \end{proof}
  \end{enumerate}
\end{enumerate}
\subsection*{Euclid's Lemma}
\begin{adjustwidth}{2em}{2em}
  \textit{Lemma. } Let $a,b,c\in\Z$. If $a\mid bc$ and $\gcd(a,b) = 1$, then $a\mid c$. 
  \begin{proof}
    If $a\mid bc$ then let $x, y \in \Z$ such that:
    \begin{align*}
      ax + by = 1
    \end{align*}
    Also let $K\in \Z$ such that $Ka = bc$
    \begin{align*}
      xac + ybc = c
    \end{align*}
    Now it gets yummier:
    \begin{align*}
      xac + yKa = c
    \end{align*}
    Look:
    \begin{align*}
      a(xc + yK) = c
    \end{align*}
    Thus $a\mid c$.
  \end{proof}
\end{adjustwidth}


\subsection*{Primes}
\subsubsection*{Prime Definition}
\begin{adjustwidth}{2em}{2em}\gspc
  \textit{Definition. }An integer $p>1$ is \ul{prime} if the only positive divisors of $p$ are 1 and $p$. 
\end{adjustwidth}\gspc
\begin{adjustwidth}{2em}{2em}\gspc
  \textit{Theorem. }Let $p\mid mn$. Then $p\mid m$ or $p\mid n$. Assume $p\nmid m$; Prove $p\mid n$.  
\end{adjustwidth}\gspc
\begin{adjustwidth}{2em}{2em}\gspc
  \textit{Theorem. }Let $p\mid m_1m_2\ldots m_k$. Then $p\mid m_i$ for at least one;
\end{adjustwidth}
\subsubsection*{Fundamental Theorem of Arithmetic}
\begin{adjustwidth}{2em}{2em}
  \textit{Theorem. }Let $n>1$ be an integer. There is a unique way to express $n$ as a product of primes, up to order 10.
\end{adjustwidth}\gspc
\subsubsection*{Infinite Primes}
\begin{adjustwidth}{2em}{2em}
  \textit{Theorem. }There are infinitely many primes
\end{adjustwidth}\gspc
\subsubsection*{Dirichlet's Theorem}
\begin{adjustwidth}{2em}{2em}
  \textit{Theorem. }Consider the sequence:
  \begin{align*}
    a, a+b, a+2b, a+3b
  \end{align*}
  If $\gcd(a,b)=1$, this sequence contains $\infty$-many primes.
\end{adjustwidth}\gspc
\pagebreak
\subsubsection*{}


\end{document}
